\documentclass[11pt]{article}

\usepackage[a4paper,margin=1in]{geometry}
\usepackage{amsmath,amssymb,bm}
\usepackage{booktabs}
\usepackage{microtype}

\title{A Linear Reaction--Diffusion Problem on the Unit Disk}
\author{}
\date{}

\newcommand{\R}{\mathbb{R}}
\newcommand{\dx}{\,\mathrm{d}x}

\begin{document}
\maketitle

\section{Domain and governing equation}

Let
\[
    \Omega=\left\{(x,y)\in\R^2:x^2+y^2<1\right\}
\]
be the unit disk, with boundary $\partial\Omega$ and outward unit normal
$\bm n$. For a final time $T=5$, the scalar field
$u:\overline{\Omega}\times[0,T]\to\R$ satisfies
\begin{equation}
    \frac{\partial u}{\partial t}
    -\nabla\cdot\bigl(a(x,y)\nabla u\bigr)
    +r(x,y)u=0
    \qquad\text{in }\Omega\times(0,T].
    \label{eq:pde}
\end{equation}
This is a linearized monodomain, or passive-cable, model. The diffusion
coefficient $a$ represents conductivity, while the nonnegative reaction
coefficient $r$ represents passive leakage or repolarization.

\section{Initial and boundary conditions}

The initial condition is an off-centre Gaussian depolarization,
\begin{equation}
    u(x,y,0)=u_0(x,y)
    =A_{\mathrm{stim}}
    \exp\!\left(
        -\frac{(x-x_0)^2+(y-y_0)^2}{2s_{\mathrm{stim}}^2}
    \right),
    \label{eq:initial}
\end{equation}
where
\[
    A_{\mathrm{stim}}=1,
    \qquad s_{\mathrm{stim}}=0.12,
    \qquad (x_0,y_0)=(0.4,0).
\]

The disk is electrically insulated, which gives the homogeneous Neumann
boundary condition
\begin{equation}
    \bigl(a(x,y)\nabla u\bigr)\cdot\bm n=0
    \qquad\text{on }\partial\Omega\times(0,T].
    \label{eq:boundary}
\end{equation}
Thus, no diffusive flux crosses the boundary.

\section{Spatial coefficients}

The conductivity is a smooth heterogeneous field containing a
low-conductivity scar:
\begin{equation}
    a(x,y)=a_{\mathrm{healthy}}
    -\bigl(a_{\mathrm{healthy}}-a_{\mathrm{scar}}\bigr)
    \exp\!\left(
        -\frac{(x-x_a)^2+(y-y_a)^2}{2w_a^2}
    \right),
    \label{eq:diffusivity}
\end{equation}
with
\[
    a_{\mathrm{healthy}}=0.1,
    \qquad a_{\mathrm{scar}}=0.01,
    \qquad (x_a,y_a)=(-0.3,0),
    \qquad w_a=0.2.
\]
At the scar centre, $a=a_{\mathrm{scar}}$, and away from the scar the
coefficient approaches $a_{\mathrm{healthy}}$. In particular,
$a(x,y)\geq a_{\mathrm{scar}}>0$, so the diffusion operator is uniformly
elliptic.

The reaction coefficient is spatially uniform:
\begin{equation}
    r(x,y)=r_0=1.
    \label{eq:reaction}
\end{equation}

\section{Continuous variational formulation}

Set $V=H^1(\Omega)$ and define the $L^2(\Omega)$ inner product by
\[
    (p,q)_\Omega=\int_\Omega pq\dx.
\]
Multiplying~\eqref{eq:pde} by an arbitrary test function $v\in V$ and
integrating the diffusion term by parts gives
\[
    (\partial_t u,v)_\Omega
    +\bigl(a\nabla u,\nabla v\bigr)_\Omega
    +(ru,v)_\Omega
    -\int_{\partial\Omega}(a\nabla u\cdot\bm n)v\,\mathrm{d}s=0.
\]
The boundary integral vanishes by~\eqref{eq:boundary}. Therefore, the weak
problem is: find $u(t)\in V$ for $t\in(0,T]$, with $u(0)=u_0$, such that
\begin{equation}
    (\partial_t u,v)_\Omega+B(u,v)=0
    \qquad\forall v\in V,
    \label{eq:weak}
\end{equation}
where
\begin{equation}
    B(w,v)
    =\int_\Omega a(x,y)\nabla w\cdot\nabla v\dx
     +\int_\Omega r(x,y)wv\dx.
    \label{eq:bilinear}
\end{equation}
The homogeneous Neumann condition is natural and therefore does not need to
be imposed explicitly in the finite-element space.

\section{Finite-element and time-discrete formulation}

Let $\mathcal{T}_h$ be a triangular mesh of $\Omega$, with target element
size $h=0.05$, and let
\[
    V_h=\left\{v_h\in C^0(\overline\Omega):
    v_h\big|_K\in\mathbb{P}_1(K)\ \text{for every }K\in\mathcal{T}_h
    \right\}
\]
be the continuous piecewise-linear finite-element space. Divide $[0,T]$
into $N=250$ equal steps, so that
\[
    \Delta t=\frac{T}{N}=0.02,
    \qquad t_n=n\Delta t.
\]

For a parameter $\theta\in[0,1]$, the fully discrete $\theta$-scheme is:
given $u_h^n\in V_h$, find $u_h^{n+1}\in V_h$ such that, for every
$v_h\in V_h$,
\begin{equation}
    (u_h^{n+1},v_h)_\Omega
    +\Delta t\,\theta B(u_h^{n+1},v_h)
    =
    (u_h^n,v_h)_\Omega
    -\Delta t\,(1-\theta)B(u_h^n,v_h).
    \label{eq:theta}
\end{equation}
The implemented choice is $\theta=1$, which gives the backward Euler
method. Choosing $\theta=\tfrac12$ instead would give the Crank--Nicolson
method. The initial finite-element function $u_h^0$ is obtained by nodal
interpolation of~\eqref{eq:initial} into $V_h$.

For $\theta=1$, equation~\eqref{eq:theta} reduces to
\begin{equation}
    (u_h^{n+1},v_h)_\Omega
    +\Delta t\,B(u_h^{n+1},v_h)
    =(u_h^n,v_h)_\Omega
    \qquad\forall v_h\in V_h.
\end{equation}
Because $a(x,y)>0$ and $r(x,y)\geq0$, the mass-plus-reaction--diffusion
system is symmetric positive definite; no nullspace treatment is required.

\end{document}
